% Generated draft: requires mathematical and source review.
\documentclass[11pt]{article}
\usepackage[margin=25mm]{geometry}
\usepackage{fontspec}
\setmainfont{Latin Modern Roman}
\usepackage{amsmath,amssymb,mathtools,unicode-math}
\setmathfont{Latin Modern Math}
\usepackage{enumitem,needspace,xurl,hyperref}
\setlength{\parindent}{0pt}
\setlength{\parskip}{5pt}
\setlength{\emergencystretch}{3em}
\newcommand{\N}{\mathbb{N}}
\newcommand{\Z}{\mathbb{Z}}
\newcommand{\Q}{\mathbb{Q}}
\newcommand{\R}{\mathbb{R}}
\newcommand{\C}{\mathbb{C}}
\newcommand{\F}{\mathbb{F}}
\newcommand{\A}{\mathbb{A}}
\newcommand{\PP}{\mathbb{P}}
\newcommand{\E}{\mathbb{E}}
\newcommand{\eps}{\varepsilon}
\newcommand{\dd}{\,\mathrm{d}}
\begin{document}
\section*{UnsolvedMath ID: 30002574}
\textbf{Source problem code:} catalogue entry
\subsection{Definitions and mathematical statement}
The source poses a modelling problem concerning cancer-cell movement through tissue containing obstacles, extracellular matrix, and fibers. A pore-passage criterion is intended to describe whether a cell can pass through a pore under the restrictions imposed by this environment. The source does not specify the geometry, admissible cell configurations, constitutive laws, or mathematical definition of successful passage.

Formulate a dynamic, time-dependent description of these restrictions and determine whether it can replace the unspecified static squeezing criterion by a condition for pore passage. The source suggests, as an example, the schematic relation
\[
F-hv=0.
\]
The quantities \(F\), \(h\), and \(v\) are not defined in the supplied record. Their scalar or vector character, units, dependence on time or configuration, and connection to traction and stiffness remain unspecified. In particular, this relation is not supplied as an established criterion, and the record does not say whether it is intended to be necessary, sufficient, or both for passage.

A precise mathematical formulation therefore remains conditional on specifying the governing model, spatial and temporal domains, initial and boundary conditions, the static criterion to be replaced, and the event that constitutes passage. No quantification over models, tissue geometries, parameters, or trajectories is provided. The supplied question cannot yet be stated as a fully specified theorem or conjecture without additional modelling assumptions.
\subsection{Short English statement}
Describe how obstacles, extracellular matrix, and fibers restrict cancer-cell motion, and formulate a time-dependent criterion for passage through a pore. The suggested equation \(F-hv=0\) is schematic: its symbols, governing model, and relationship to successful passage require clarification.
\subsection{Sources}
\begin{itemize}
\item Local UnsolvedMath registry, \texttt{data/unsolvedmath/problems.json}, permanent ID \(30002574\), source problem code \texttt{catalogue entry}, title ``Dynamic Constraints on Cancer-Cell Passage Through Tissue.'' The supplied record belongs to the Oberwolfach Reports Open Problems collection; no individual proposer is identified.
\item Supplied source citation: \emph{Mini-Workshop: Mathematical Models for Cancer Cell Migration (2015)}, Oberwolfach Reports, DOI 10.4171/owr/2014/19, pp. 1075--1109. \url{https://doi.org/10.4171/owr/2014/19}. The record describes Oberwolfach Reports as documenting mathematical workshops and their open-problem discussions. Its proposed year and year tag are 2014, whereas the supplied citation names 2015; this discrepancy is unresolved.
\item Supplied literature source, labelled ``Micropore immersed-boundary model'': \url{https://doi.org/10.1177/0954406217730850}.
\item Supplied literature source, labelled ``Constricted-channel cancer-cell model'': \url{https://doi.org/10.1007/s10237-023-01705-6}.
\item Literature assessment dated exactly 2026-08-21: the supplied status is ``partially solved.'' The supplied assessment states that time-dependent computational and experimental models now relate pore size, pressure, membrane rheology, nuclear stiffness, and transit time, but that a general dynamic squeezing criterion in terms of traction and stiffness has not been derived or validated. This is supplied dated curation derived from the OWR record, not an independently verified literature assessment. The background explicitly requests independent verification of time-sensitive status and cited claims.
\item The supplied statement status is \texttt{corrected\_verified}, with the explanation that mathematical content was preserved while notation, typography, or omitted definitions were repaired for publication. This is a supplied editorial designation, not independent verification of the mathematical formulation or literature status.
\end{itemize}
\end{document}
